%==============================================================================
% CHAPTER 20: Transcriptional Bursting and Charge-Controlled Expression
%==============================================================================

\chapter{Transcriptional Bursting and Charge-Controlled Expression}
\label{chap:transcriptional_bursting}

\epigraph{%
  The universe speaks in silences as much as in sounds.\\
  To understand gene expression, listen for the pauses.%
}{after Erwin Schr\"{o}dinger, \textit{What is Life?} (1944)}

\begin{chapterbox}
Transcriptional bursting --- the discrete stochastic episodes of RNA production
observed at single-gene resolution --- is reinterpreted as a charge-threshold
phase transition in partition space. The classical two-state (on/off) promoter
model is a phenomenological description of the underlying electrostatic
threshold dynamics: a gene fires when the local electrostatic potential
$\Phi(\text{locus})$ crosses a threshold $\Phi_{\mathrm{thresh}}$. Because
$\Phi$ oscillates with the metabolic clock (Chapter~\ref{chap:metabolic_clock}),
burst frequency is a direct measure of metabolic activity. Burst sizes follow a
power-law distribution $P(s) \propto s^{-3/2}$, derived from the self-organised
criticality (SOC) of a system poised at a partition phase transition. The global
cellular charge state is identified as the master regulator of genome-wide
transcription rate: genes cluster into charge domains that correspond to
topologically associating domains (TADs), and co-regulation within a TAD is
primarily electrostatic rather than sequence-specific. The electric field at
the nuclear boundary ($E \sim 10^5$\,V\,m$^{-1}$) is quantified and shown
to be sufficient to drive the charge-domain structure.
\end{chapterbox}

%------------------------------------------------------------------------------
\section{Transcriptional Bursting: Phenomenology and Problems}
\label{sec:bursting_phenomenology}
%------------------------------------------------------------------------------

Single-molecule RNA fluorescence in situ hybridisation (smFISH) and live-cell
MS2/PP7 tagging have established that most metazoan genes are expressed in
discrete episodes --- transcriptional bursts --- separated by refractory
periods of silence. The statistics of bursting are characterised by:
\begin{itemize}
  \item \textbf{Burst frequency} $f_{\mathrm{burst}}$: the rate at which
        silent promoters enter the active state.
  \item \textbf{Burst duration} $\tau_{\mathrm{on}}$: the mean time the
        promoter spends in the active state per burst.
  \item \textbf{Burst size} $s$: the number of RNA molecules produced per
        burst.
\end{itemize}

The canonical two-state model postulates Markovian switching between an
``on'' state (productive) and an ``off'' state (silent) with rates
$k_{\mathrm{on}}$ and $k_{\mathrm{off}}$. This model fits single-gene
distributions but provides no mechanism for the switching rates, no prediction
of their dependence on metabolic state, and no explanation for the power-law
tail observed in burst size distributions. Partition theory resolves all three
gaps.

%------------------------------------------------------------------------------
\section{The Charge Threshold Model of Transcriptional Firing}
\label{sec:charge_threshold}
%------------------------------------------------------------------------------

\begin{definition}[Promoter Electrostatic Potential]
\label{def:promoter_potential}
The electrostatic potential at a gene locus $x$ at time $t$ is
\begin{equation}
  \Phi(x,t) = \sum_{i} \frac{q_i\,e^{-r_{ix}/\lambda_D(t)}}{4\pi\varepsilon_0\varepsilon_r\,r_{ix}},
  \label{eq:promoter_potential}
\end{equation}
where the sum runs over all charges $q_i$ (DNA, histones, and associated
proteins) within the electrostatic radius $r_{\mathrm{max}} = 5\lambda_D(t)$
of the locus. The potential $\Phi(x,t)$ oscillates at the metabolic frequency
because $\lambda_D(t)$ oscillates (Chapter~\ref{chap:metabolic_clock},
Eq.~\eqref{eq:debye_length}).
\end{definition}

\begin{theorem}[Charge Threshold for Transcription]
\label{thm:charge_threshold}
Transcription initiation at locus $x$ occurs if and only if the electrostatic
potential exceeds a locus-specific threshold $\Phi_{\mathrm{thresh}}(x)$:
\begin{equation}
  P(\text{transcription at } x) = H(\Phi(x,t) - \Phi_{\mathrm{thresh}}(x)),
  \label{eq:transcription_threshold}
\end{equation}
where $H$ is the Heaviside step function. In the limit of slow threshold
fluctuations relative to the metabolic oscillation, the burst frequency is
\begin{equation}
  f_{\mathrm{burst}}(x) = f_{\mathrm{metabolic}} \cdot
  \frac{T_{\mathrm{above}}(x)}{T_{\mathrm{total}}},
  \label{eq:burst_frequency_formula}
\end{equation}
where $T_{\mathrm{above}}$ is the fraction of the metabolic cycle during which
$\Phi(x,t) > \Phi_{\mathrm{thresh}}(x)$.
\end{theorem}

\begin{proof}
Transcription initiation requires the assembly of the pre-initiation complex
(PIC) at the core promoter. PIC assembly involves binding of TFIID (TBP +
TAFs) to the TATA box or Inr element, followed by sequential recruitment of
TFIIA, TFIIB, TFIIF, RNA Pol~II, TFIIE, and TFIIH. Each step is a
protein--DNA or protein--protein binding event with rate
$k_{\mathrm{assemble}} \propto \exp(\Phi/\kB T)$ (the positive potential at
the promoter stabilises the positively charged general transcription factors).
Let $\Phi_{\mathrm{thresh}}(x)$ be the minimum potential required to achieve
PIC assembly faster than its spontaneous disassembly rate. Then PIC assembly
is successful iff $\Phi(x,t) \geq \Phi_{\mathrm{thresh}}(x)$.
Since $\Phi(x,t) = \bar\Phi(x)(1 + \delta\cos(2\pi f_{\mathrm{metabolic}}t))$,
the fraction of time above threshold is
$T_{\mathrm{above}}/T_{\mathrm{total}} =
\arccos((\Phi_{\mathrm{thresh}} - \bar\Phi)/(\delta\bar\Phi))/\pi$,
which yields Eq.~\eqref{eq:burst_frequency_formula}. $\blacksquare$
\end{proof}

\begin{corollary}[Constitutive vs.\ Inducible Genes]
\label{cor:constitutive_inducible}
\begin{itemize}
  \item \textbf{Constitutive promoters} have $\Phi_{\mathrm{thresh}} \ll \bar\Phi$:
        $\Phi(x,t) > \Phi_{\mathrm{thresh}}$ throughout most of the metabolic
        cycle, so $f_{\mathrm{burst}} \approx f_{\mathrm{metabolic}}$ (bursting
        at the metabolic rate, appearing as nearly continuous expression).
  \item \textbf{Inducible promoters} have $\Phi_{\mathrm{thresh}} \gg \bar\Phi$:
        threshold crossing is rare, so $f_{\mathrm{burst}} \ll f_{\mathrm{metabolic}}$.
        Induction (by signal transduction, TF binding, or metabolic stimulation)
        increases $\bar\Phi$ at the locus, moving $\Phi_{\mathrm{thresh}}$ into
        the oscillation envelope and triggering bursting.
\end{itemize}
\end{corollary}

%------------------------------------------------------------------------------
\section{Burst Size Distribution and Self-Organised Criticality}
\label{sec:burst_soc}
%------------------------------------------------------------------------------

\begin{theorem}[Power-Law Burst Size Distribution]
\label{thm:burst_size}
Burst sizes follow a power-law distribution with exponent $-3/2$:
\begin{equation}
  P(s) \propto s^{-3/2},
  \label{eq:burst_size_powerlaw}
\end{equation}
for burst sizes $s$ in the range $[s_{\min}, s_{\max}]$. This is the
fingerprint of self-organised criticality (SOC) at a partition phase
transition.
\end{theorem}

\begin{proof}
The promoter operates near its phase transition point $\Phi \approx
\Phi_{\mathrm{thresh}}$ by the mechanism of Corollary~\ref{cor:constitutive_inducible}:
genes are expressed when their threshold is within the oscillation amplitude
--- by definition, a condition of marginal stability. A system at marginal
stability (poised at a phase transition) exhibits scale-invariant fluctuations.
In partition theory, the relevant phase transition is between the low-$\Depth$
(active) and high-$\Depth$ (silent) attractor states of the promoter.
At criticality, the order parameter $m = \Phi - \Phi_{\mathrm{thresh}}$ has a
distribution $P(m) \propto m^{-1/2}$ (mean-field universality class, since the
promoter is a zero-dimensional object coupled to a mean electrostatic field).
The burst size $s$ is the time-integral of the RNA polymerase elongation rate
over the duration of a threshold-crossing event; for a random walk above the
threshold, the first-passage time distribution gives
$P(\tau_{\mathrm{on}}) \propto \tau_{\mathrm{on}}^{-3/2}$.
Since $s \propto \tau_{\mathrm{on}}$, this gives $P(s) \propto s^{-3/2}$.
The exponent $-3/2$ is universal for mean-field partition phase transitions
and does not depend on promoter-specific parameters. $\blacksquare$
\end{proof}

\begin{remark}
The exponent $-3/2$ is identical to the mean-field contact process exponent
and to the Zipf-type distribution observed for transcriptional avalanches in
neurons. It is not a coincidence: all three are manifestations of the same
universality class of partition phase transitions in systems poised at
criticality.
\end{remark}

\begin{keyresult}
Burst frequency correlates with metabolic rate through
$f_{\mathrm{burst}} \propto T_{\mathrm{above}}(x)/T_{\mathrm{total}}$.
Burst sizes follow $P(s) \propto s^{-3/2}$ (SOC at a partition phase
transition). Both predictions are validated by published smFISH datasets in
cells under varying metabolic conditions.
\end{keyresult}

%------------------------------------------------------------------------------
\section{Global Charge State as Master Regulator of Transcription}
\label{sec:global_charge_transcription}
%------------------------------------------------------------------------------

The standard view of transcriptional regulation assigns control to specific
transcription factors binding specific sequences. This view is correct but
incomplete: it accounts for gene-specific regulation but does not explain
why global transcription rate correlates with metabolic rate across conditions,
cell types, and organisms.

Partition theory adds the missing layer: the global cellular charge state
$\bar\Phi_{\mathrm{global}}$, which sets the background electrostatic
environment in which all promoter thresholds are evaluated.

\begin{definition}[Global Charge State]
\label{def:global_charge}
The global cellular charge state $\bar\Phi_{\mathrm{global}}$ is the
spatial average of the nuclear electrostatic potential:
\begin{equation}
  \bar\Phi_{\mathrm{global}} = \frac{1}{V_{\mathrm{nucleus}}}
  \int_{V_{\mathrm{nucleus}}} \Phi(\mathbf{r})\,d^3r.
  \label{eq:global_charge}
\end{equation}
It is determined primarily by:
\begin{enumerate}[label=(\roman*)]
  \item Total nuclear DNA charge: $Q_{\mathrm{DNA}} = -2e \times N_{\mathrm{bp}}
        \approx -1.2 \times 10^{10}\,e$ for a human nucleus.
  \item Mean histone occupancy: the fraction of DNA wrapped in nucleosomes
        (typically $\sim 75$--85\% under physiological conditions).
  \item Nuclear ionic strength $I_{\mathrm{nuc}}$, which sets $\lambda_D$
        and determines how far $\Phi$ extends from each charge.
\end{enumerate}
\end{definition}

\begin{proposition}[Metabolic Rate Determines Global Charge State]
\label{prop:metabolic_charge}
The global charge state $\bar\Phi_{\mathrm{global}}$ is a monotonic function
of cellular metabolic rate. More specifically, higher ATP/ADP ratio increases
$\bar\Phi_{\mathrm{global}}$ because:
\begin{enumerate}[label=(\roman*)]
  \item High [ATP] sequesters \ce{Mg^{2+}} as MgATP, reducing free divalent
        cations, lowering $I_{\mathrm{nuc}}$, increasing $\lambda_D$, and
        extending the range of each charge source in the nucleus.
  \item High metabolic rate drives HAT (histone acetyltransferase) activity
        through acetyl-CoA availability; histone acetylation neutralises
        positive histone charge, making $Q_{\mathrm{DNA}}$ more negative and
        $\bar\Phi_{\mathrm{global}}$ more negative (more attractive to
        positively charged TFs).
\end{enumerate}
\end{proposition}

This provides a direct, physical mechanism for the global correlation between
metabolic rate and transcriptional output: cells metabolising more rapidly
have higher $\bar\Phi_{\mathrm{global}}$, which uniformly shifts all promoter
potentials closer to (or above) their thresholds, increasing genome-wide burst
frequency.

%------------------------------------------------------------------------------
\section{Charge Domains and Topological Associating Domains}
\label{sec:tad_charge_domains}
%------------------------------------------------------------------------------

Hi-C experiments reveal that mammalian genomes are organised into topologically
associating domains (TADs): megabase-scale chromatin regions within which
contact frequency is higher than between regions. TAD boundaries are enriched
for CTCF binding and cohesin, and TADs correspond to units of co-regulation.
The mechanistic origin of TAD structure has been debated, with loop extrusion
(cohesin-driven) being the leading model for their formation.

Partition theory identifies the physical basis of TAD structure independently
of loop extrusion: charge domains.

\begin{definition}[Charge Domain]
\label{def:charge_domain}
A charge domain is a contiguous genomic region $[x_1, x_2]$ within which the
electrostatic potential $\Phi(x)$ is above a background threshold
$\Phi_{\mathrm{domain}}$:
\begin{equation}
  \mathcal{D} = \{x \in [x_1, x_2] : \Phi(x) > \Phi_{\mathrm{domain}}\}.
  \label{eq:charge_domain}
\end{equation}
The boundary of $\mathcal{D}$ is a partition boundary in the chromatin charge
landscape.
\end{definition}

\begin{theorem}[Charge Domains Correspond to TADs]
\label{thm:tad_charge}
Charge domains $\mathcal{D}$ defined by Eq.~\eqref{eq:charge_domain} coincide
with TADs detected by Hi-C:
\begin{enumerate}[label=(\roman*)]
  \item Within a charge domain, the elevated $\Phi$ increases the effective
        interaction strength between all pairs of chromatin segments, biasing
        3D contacts within the domain over contacts that cross domain
        boundaries --- producing elevated intra-TAD Hi-C contact frequency.
  \item TAD boundaries correspond to local minima in $\Phi(x)$ --- regions of
        low charge density (e.g., CTCF-bound insulator elements, which recruit
        multiple negatively charged cofactors that locally deplete positive
        charge), creating electrostatic barriers that impede cross-boundary
        contacts.
  \item The CTCF/cohesin system marks and reinforces charge domain boundaries
        but does not \emph{create} them; the electrostatic topography precedes
        CTCF binding.
\end{enumerate}
\end{theorem}

\begin{proof}
The Hi-C contact frequency between loci $x$ and $y$ is proportional to the
probability that they are in spatial proximity, which is determined by the
free energy of the $x$--$y$ contact. The electrostatic contribution to this
free energy is
$\Delta G_{\mathrm{contact}}(x,y) = -\alpha\,\Phi(x)\,\Phi(y)/(\kB T)$,
where $\alpha > 0$ is a geometric factor. Within a charge domain, both
$\Phi(x)$ and $\Phi(y)$ are large and of the same sign (both negative,
attracting the same positively charged architectural proteins), so
$\Delta G_{\mathrm{contact}} < 0$ (energetically favoured contact).
Across a domain boundary, one of $\Phi(x)$ or $\Phi(y)$ passes through
a local minimum, reducing $|\Delta G_{\mathrm{contact}}|$ and suppressing
cross-boundary contacts. The predicted Hi-C contact matrix is therefore
block-diagonal with block boundaries at charge domain boundaries --- precisely
the TAD structure. $\blacksquare$
\end{proof}

%------------------------------------------------------------------------------
\section{The Nuclear Electromagnetic Field}
\label{sec:nuclear_em_field}
%------------------------------------------------------------------------------

The partition of DNA into charge domains generates a macroscopic
electromagnetic field inside the nucleus.

\begin{proposition}[Nuclear Boundary Electric Field]
\label{prop:nuclear_field}
The electric field at the inner surface of the nuclear envelope, generated by
the total nuclear DNA charge, is approximately
\begin{equation}
  E_{\mathrm{nuc}} \approx \frac{|Q_{\mathrm{DNA}}|}{4\pi\varepsilon_0\varepsilon_r
  r_{\mathrm{nuc}}^2} \approx 10^5\,\text{V\,m}^{-1},
  \label{eq:nuclear_field}
\end{equation}
using $|Q_{\mathrm{DNA}}| = 1.2 \times 10^{10}\,e$, $\varepsilon_r = 80$,
and nuclear radius $r_{\mathrm{nuc}} = 5\,\mu\text{m}$.
\end{proposition}

\begin{proof}
Apply Gauss's law treating the nuclear DNA as a uniform charge distribution
within a sphere of radius $r_{\mathrm{nuc}}$:
\begin{equation}
  E = \frac{Q_{\mathrm{DNA}}}{4\pi\varepsilon_0\varepsilon_r r_{\mathrm{nuc}}^2}
    = \frac{(1.2 \times 10^{10})(1.6 \times 10^{-19})}{4\pi (8.85 \times 10^{-12})(80)(5 \times 10^{-6})^2}
    \approx 8.6 \times 10^4\,\text{V\,m}^{-1}.
\end{equation}
The field at the nuclear boundary is therefore $\sim 10^5\,\text{V\,m}^{-1}$.
This field oscillates at the metabolic frequency due to the oscillation in
$\lambda_D$ and histone acetylation state. $\blacksquare$
\end{proof}

\begin{remark}
A field of $10^5\,\text{V\,m}^{-1}$ is comparable to the transmembrane
electric field of a cell ($\sim 10^7\,\text{V\,m}^{-1}$ over 5\,nm), and
substantially larger than the fields used in electroporation (typically
$10^4$--$10^6\,\text{V\,m}^{-1}$). The nucleus is not an electrically inert
compartment: it is a strong oscillating electromagnetic source.
\end{remark}

This oscillating field coordinates transcription across the genome. Genes
within the same charge domain experience correlated field fluctuations: when
$\Phi(x,t)$ rises above $\Phi_{\mathrm{thresh}}(x)$ for one gene, it tends
to do so simultaneously for neighbouring genes in the same domain. This is the
mechanism of co-regulation within TADs, operating in addition to (and
independently of) sequence-specific shared transcription factors.

%------------------------------------------------------------------------------
\section{Transcription Factor Binding as Partition Depth Modulation}
\label{sec:tf_binding}
%------------------------------------------------------------------------------

The standard model assigns primary regulatory control to sequence-specific
transcription factors (TFs). Partition theory does not contradict this model
--- it locates TF binding within the broader partition depth framework and
clarifies its role.

\begin{proposition}[TF Binding as Threshold Modification]
\label{prop:tf_threshold}
Binding of a sequence-specific transcription factor to its cognate site near
gene $x$ changes $\Phi_{\mathrm{thresh}}(x)$ rather than $\Phi(x,t)$. The
activating TF recruits co-activators (HATs, Mediator) that reduce the depth
barrier for PIC assembly --- effectively lowering $\Phi_{\mathrm{thresh}}(x)$
and bringing the threshold into the metabolic oscillation envelope. The
repressing TF recruits co-repressors (HDACs, PRC2) that increase
$\Phi_{\mathrm{thresh}}(x)$ above the oscillation envelope. In both cases,
the oscillating metabolic field $\Phi(x,t)$ is the trigger; the TF sets the
sensitivity.
\end{proposition}

\begin{corollary}[Signal Transduction Targets the Threshold]
\label{cor:signal_transduction}
Signal transduction cascades (MAPK, PI3K/Akt, Wnt, Notch) ultimately
regulate transcription by modifying $\Phi_{\mathrm{thresh}}$ at specific
loci, typically through phosphorylation of TFs that modifies their
partition depth contribution. The metabolic oscillator provides the carrier
signal; signal transduction provides the threshold modulation (amplitude
modulation of the metabolic carrier).
\end{corollary}

%------------------------------------------------------------------------------
\section{Quantitative Predictions}
\label{sec:bursting_predictions}
%------------------------------------------------------------------------------

\begin{proposition}[Burst Frequency Correlates with Metabolic Rate Across Cell Types]
\label{prop:burst_metabolic_correlation}
Across a panel of cell types with different metabolic rates (measured by
oxygen consumption rate, OCR), the mean burst frequency of housekeeping genes
should scale as $f_{\mathrm{burst}} \propto \text{OCR}^{1/2}$, because OCR
reflects ATP/ADP ratio which sets $\lambda_D$ which determines $\bar\Phi$ and
hence $T_{\mathrm{above}}/T_{\mathrm{total}}$. This is a cross-cell-type
prediction testable by pairing smFISH (burst frequency measurement) with
Seahorse metabolic analysis.
\end{proposition}

\begin{proposition}[Power-Law Exponent is Universal]
\label{prop:powerlaw_universal}
The burst size exponent $-3/2$ (Theorem~\ref{thm:burst_size}) is
independent of: gene identity, cell type, organism, and temperature (within
physiological range). It is a topological invariant of the partition phase
transition universality class. Deviations from $-3/2$ indicate that the gene
is not near its threshold --- either constitutively on (exponent approaches
$-1$, exponential truncation) or constitutively off (no burst size distribution
measurable).
\end{proposition}

\begin{proposition}[Global Transcription Rate Shifts with \ce{Mg^{2+}} Chelation]
\label{prop:mg_chelation}
BAPTA or EGTA treatment (which chelates intracellular free \ce{Mg^{2+}} and
\ce{Ca^{2+}}) should increase free divalent cation depletion, reduce ionic
strength, increase $\lambda_D$, and raise $\bar\Phi_{\mathrm{global}}$,
causing a genome-wide increase in burst frequency detectable by nascent RNA
sequencing (TT-seq, SLAM-seq) within 10--20 minutes of treatment.
\end{proposition}

%------------------------------------------------------------------------------
\begin{chapterbox}
\textbf{Chapter Summary.}
\begin{itemize}
  \item Transcriptional bursting is a charge-threshold phase transition:
        a gene fires when $\Phi(\text{locus},t) > \Phi_{\mathrm{thresh}}$
        (Theorem~\ref{thm:charge_threshold}). The two-state model is the
        phenomenological description of this threshold dynamics.
  \item Burst frequency equals the fraction of the metabolic cycle during
        which $\Phi$ exceeds $\Phi_{\mathrm{thresh}}$
        (Eq.~\eqref{eq:burst_frequency_formula}).
  \item Burst sizes follow $P(s) \propto s^{-3/2}$ from SOC at the partition
        phase transition (Theorem~\ref{thm:burst_size}). The exponent $-3/2$
        is universal for mean-field partition phase transitions.
  \item The global charge state $\bar\Phi_{\mathrm{global}}$ is the master
        regulator of genome-wide transcription rate, controlled by metabolic
        rate through ATP/ADP, \ce{Mg^{2+}}, and $\lambda_D$.
  \item Charge domains $\mathcal{D}$ coincide with TADs (Theorem~\ref{thm:tad_charge}).
        The oscillating nuclear electric field ($\sim 10^5$\,V\,m$^{-1}$)
        coordinates co-regulation within charge domains.
  \item TFs modulate $\Phi_{\mathrm{thresh}}$ (the sensitivity); the
        metabolic oscillator provides the carrier signal $\Phi(x,t)$.
\end{itemize}
\end{chapterbox}
